\section*{Matched consent comparison}
The primary consent analysis reuses the nine fitted predictors and their original test splits. Model parameters, training data, labels, patient membership and validation-selected decision thresholds remain unchanged. The public full-input condition is the reference; it is not a claim that clinical data may be used without permission. The intervention applies only to inference-time availability. Consent to training, removal of training records, machine unlearning, actual preference collection and security against an attacker are outside this experiment.

There are eight prespecified conditions: full input; independent laboratory and vital-sign denial with probability 0.1, 0.3 or 0.5 per group; no laboratory inputs; no vital inputs; metadata only; and denial with probability 0.3 per group beginning at zero-based hour 12. Metadata remain available in every condition. Each patient's two random permission draws are fixed across models and rates, so higher denial rates form nested restrictions for that patient. The draws are keyed by patient, site and split seed and do not use outcomes. Probabilities concern each feature group; the probability of at least one denied group is $1-(1-p)^2$ at common rate $p$.

A denied value and its running mean become missing, while its observation-count features become zero. On withdrawal, denied historical features are cleared at the specified hour. A GRU restarts its recurrent state at that boundary, preventing its prior mixed state from carrying the denied information forward. This conservative reset also discards permitted historical context. Model imputation and normalization retain the transformations learned during the original training; a restricted model is not retrained or silently recalibrated.

For both first-alert and repeated-hourly policies, the threshold maximizes challenge utility on the original validation population and is then reused in every consent condition. Thus a paired effect includes the consequences of restricted inputs and the specified state reset under a fixed policy. It does not include a compensatory threshold change. The earlier scenario-adaptive consent analysis remains a separate secondary analysis and is not pooled with these paired effects.

Discrimination uses the same pre-onset rows in both conditions. We report AUROC, average precision and Brier score. Two hundred paired whole-patient bootstrap replicates estimate pointwise percentile intervals for differences in AUROC and Brier score; average precision is reported as a point estimate. Five hundred paired patient resamples estimate differences in conditional incident missed risk, reviews per patient, normalized challenge utility and false reviews per benchmark-negative patient. Each replicate uses identical patient multiplicities in the reference and restricted conditions, preserving the covariance of the comparison. Conditional missed risk recomputes its positive-patient denominator within each resample. Intervals condition on the trained model, test split and simulated permission assignments; they do not include retraining or real preference uncertainty.

We additionally count permitted physiological source-hours and the fraction of originally available last-value feature cells that remain available. These are information-availability proxies, not estimates of leaked information, re-identification probability or privacy protection. All scenarios retain the same outcome labels for evaluation. The experiment does not suppress treatment or simulate a change in disease caused by withholding data from the model.

The protocol specifies no clinically justified non-inferiority margin. Consequently, an interval containing zero is not interpreted as proof of equal performance, and an arbitrarily chosen tolerance is not presented as an ethical limit. Results are reported as signed changes with uncertainty for each model, hospital and split. Across-seed means and ranges are descriptive because the patient populations overlap.


We evaluate both a repeated-hourly policy and a first-only policy. At threshold $\tau$, the repeated policy emits a review at every hour with $s_{it}\ge\tau$. Its missed-incident loss is one if no such hour lies in the six-hour warning window. The first-only policy emits at the first crossing and succeeds only if that crossing lies in the window; an earlier crossing consumes the review. Reviews are counted over the retrospective monitoring interval specified above. The repeated policy can exceed one review per patient, whereas the first-only policy cannot. These operational definitions prevent an improvement in ranking from being mistaken for unchanged incident detection.

The normalized challenge utility is
\[
 U=\frac{U_{\rm observed}-U_{\rm inaction}}
 {U_{\rm optimal}-U_{\rm inaction}},
\]
using the official evaluator \citep{reyna2020}. Repeated policies use the hourly threshold state. First-only policies latch a positive prediction state after their initial crossing; that persistent prediction is not counted as an additional review. Utility covers the full observed stay and uses the challenge's own timing rewards, so it differs from the strict incident-window and workload endpoints. False reviews per negative patient use only patients whose entire label sequence is negative.

AUROC measures ordering of positive and negative pre-onset hours; average precision summarizes the precision--recall curve; Brier score averages squared probability error. Whole-patient resampling preserves within-patient temporal dependence \citep{efron1979}. Because of class weighting and rare positive hours, a lower Brier score need not imply improved event detection or reliable calibration. The benchmark therefore releases all these endpoints and does not collapse them into one score.

The full-input reproduction check verifies every saved predictor against its original prediction array before any mask is applied. The resulting reference-versus-reference differences and their bootstrap intervals are exactly zero. Released patient-level policy contributions, model hashes, split assignments and permission configuration make the paired analysis independently inspectable. The benchmark protocol also separates future training-adapted or policy-adapted entries from these frozen full-input baselines.
